\section{Impulsión y medición del tratamiento de bancos: ENV-28}
\label{sec:ambiente-env28}
\textbf{Selección de impulsión.} Se seleccionan ocho Systemair K~315L~EC, artículo \texttt{449571}, 230 V monofásicos, 50/60 Hz, distribuidos en dos motores en serie por cada una de las cuatro rutas de tratamiento. Una ruta por banco opera y otra permanece de reserva. La ficha de esta variante declara 331 W y 1,459 A máximos por motor, IP54, diámetro de conexión 315 mm, envolvente de diámetro 408 mm, longitud 225 mm y aire transportado hasta 60\,$^{\circ}$C. No se trasladan datos del antiguo artículo 2585. La ficha presenta un punto numérico de 883,98 m$^3$/h a 414,85 Pa, con densidad 1,204 kg/m$^3$, 269,9 W y 1,24 A; ese punto no se convierte en garantía de 900 m$^3$/h \parencite[pp.~1--2]{lote28-k315}.

La curva gráfica a 900 m$^3$/h se lee aproximadamente en 400--410 Pa por motor. La suma ideal de dos curvas en serie da alrededor de 800--820 Pa a igual densidad; es cálculo de proyecto, pendiente de contraste de interacción, montaje y curva del conjunto. Se fija un presupuesto preliminar de 600 Pa por ruta: serpentín húmedo 150, filtro sucio 100, calefactor 25, antirretorno 30, iris 100, conductos/transiciones 80, terminales 40 y reserva 75 Pa. Las cotas de serpentín, filtro, calefactor, conductos y terminales son condiciones de selección, no pérdidas verificadas de equipos ya elegidos. El antirretorno tiene lectura gráfica cercana a 24 Pa a 900 m$^3$/h; se reserva 30 Pa. El iris requiere selección de apertura y coeficiente antes de validar su pérdida. Por tanto, 600 Pa no constituye cálculo cerrado ni habilita operación del tratamiento completo. Deben verificarse densidad local, curvas de cada componente, caudal útil y estabilidad de la pareja, con filtro sucio y conmutación.

Se colocan los motores después del recalentamiento, en aire seco tratado, antes de la medición y del plenum final; no se instalan en extracción de H$_2$. Para mantener las holguras del calefactor se usa una estación propia fuera de los 2,2 m del plenum ENV-25. Se reserva por ruta una estación recta de 2,4 por 0,8 por 0,8 m: dos cuerpos de 225 mm, separación provisional de 315 mm, antirretorno de 125 mm, iris y rectos de medida. La separación entre motores es reserva de proyecto pendiente del manual aplicable; no es distancia de fabricante verificada. La ruta se retira completa al fallar cualquiera de sus dos motores; no se presume 900 m$^3$/h a través del motor detenido. La nueva estación requiere integrar huella, accesos y conductos en el plano ejecutivo, sin afirmar que cabe en la reserva anterior de UTA de 3 por 2 m.

\textbf{Antirretorno y medición.} Cada ruta incorpora un Systemair RSK~315, artículo \texttt{5604}, de dos hojas con resorte, diámetro 314 mm y longitud 125 mm, y un iris SPI~315~C, artículo \texttt{6757}, diámetro interior 314 mm y exterior 406 mm. El SPI tiene tomas fijas y estanqueidad de conducto clase C; la ficha pide rectos de un diámetro respecto de codos y tres diámetros antes de una derivación T o terminal de impulsión. Se reserva como mínimo 945 mm rectos aguas abajo hacia el terminal; el plano debe contrastar también perturbaciones aguas arriba. El RSK evita retorno aerodinámico ordinario, pero no se atribuye estanqueidad a H$_2$, clasificación Ex ni función cortafuego a esa selección. La clase C del iris tampoco acredita aislamiento gaseoso del banco \parencite[pp.~1--2]{lote28-rsk}\parencite[pp.~1--3]{lote28-spi}.

Se seleccionan cuatro HK Instruments DPT-Flow-1000-D, código \texttt{102.001.012}, sin autozero, alimentados a 24 V y configurados en 4--20 mA para caudal de impulsión. Consumo de catálogo inferior a 1,2 W por unidad en salida de corriente: se reservan 4,8 W y 0,2 A totales a 24 V. Se realiza cero manual al comisionar y se incorpora su comprobación al mantenimiento. Se excluye la variante AZ, cuyo autozero congela la salida cuatro segundos cada diez minutos. La respuesta seleccionable mínima es un segundo, no instantánea. A 100 Pa la especificación $1\,\%+2$ Pa da hasta 3 Pa; con $Q=K\sqrt{\Delta p}$ ello representa aproximadamente 1,5\,\% de caudal por el transmisor, sin incluir incertidumbre del iris, instalación o calibración. No se inventa $K$ ni apertura: se registran del diagrama/manual del SPI y se contrastan mediante recorrido de velocidades antes de aceptar el caudal. Las tomas y transmisores permanecen en aire limpio no condensante, dentro del rango ambiente de $-20$ a $+50\,^{\circ}$C, sin comunicación neumática con extracción de H$_2$. La señal se asigna al AI WAGO 750-455 ordinario seleccionado en ENV-34, una entrada por DPT; el transmisor carece de relé de permiso. El diagnóstico analógico seleccionado no constituye función de seguridad. La selección y cableado de la cadena independiente, auxiliares de estado y contactos de inhibición deben completar la habilitación \parencite[p.~2]{lote28-dpt-ficha}\parencite{lote28-dpt-manual}.

\textbf{Demanda y calor propios.} Cuatro motores activos reservan 1,324 kW y $4(230)(1,459)/1000=1,34228$ kVA. En el solape de ocho motores son 2,648 kW y 2,68456 kVA; no se suma solape de calefactores. Cada banco activo añade 662 W de motores. Si toda esa potencia termina en la corriente tratada, el recalentamiento de proceso ENV-25 disminuye de 4,144756 a 3,482756 kW por banco; se conserva la reserva máxima eléctrica de 18,1 kW de calefactores/control. No se resta dos veces calor de motor ni se reduce por este cálculo la carga de frío del serpentín. Los 22,8604 kW por banco corresponden al punto de rocío 8,5 grados C; con el presupuesto instrumental de ENV-34 no acreditan el criterio de máximo real de 8,5 grados C y deben recalcularse para la condición más exigente. Los valores de recalentamiento aquí calculados describen ese punto anterior y no se presentan como demanda final del nuevo estado. Las fases, conductores, protecciones, fuente protegida y autonomía deben recoger la demanda nueva y el estado de solape; no se deduce corriente de arranque del valor nominal.

\textbf{Frontera de habilitación.} El caudal de 900 m$^3$/h por banco queda como condición del tratamiento desarrollado, sin habilitar la recarga agregada de 60 A. Con generación de H$_2$ de 7,5978 m$^3$/h, la concentración ideal de equilibrio resulta $7,5978/900=0,008442$, o 0,8442\,\% en volumen: supera incluso el umbral crítico de 0,8\,\%. Para mantener equilibrio inferior a 0,4\,\% sin margen se requieren más de 1.899,45 m$^3$/h; esa ampliación modifica frío, ductos, motores, energía y geometría y no se adopta por simple cambio de consigna. A 900 m$^3$/h el límite matemático proporcional para 0,4\,\% es 28,4293 A agregados; tampoco constituye corriente autorizada sin márgenes, transitorio, incertidumbre y ajuste verificable del rectificador. La medida de impulsión no demuestra distribución de extracción ni evita acumulación local. Se conserva inhibición de recarga hasta cerrar caudal útil de extracción, margen y transitorio H$_2$, clasificación y aislamiento de rutas, selección húmeda/hidráulica, fuente Carrier corregida y cadena de seguridad. RT-06.03, RT-06.09 y RT-06.13 permanecen pendientes de esas dependencias y de la integración completa.

Cada motor recibe una AO WAGO 750-559 individual según ENV-34; no se agrupan dos entradas EC de impedancia desconocida. Antes de aceptar mando se debe verificar carga de cada entrada contra el mínimo de 5 k$\Omega$ del módulo o seleccionar interfaz compatible, además de sentido, rango y ajuste PI de caudal. Los cuatro DMT132 se incorporan al aire limpio después de recalentamiento/ventiladores, sin comunicar con retorno H$_2$. La reserva de montaje no acredita cabida completa: los tramos fríos anteriores al calefactor expuestos al exterior adoptan al menos 100 mm de aislamiento conforme INT-34; 25 mm se limita al interior tratado de ENV-25.
